% !TeX program = pdflatex
\documentclass[11pt,a4paper]{amsart}

\newcommand{\notelanguage}{finnish}
% Shared preamble for course notes and exercise sets.

\usepackage[T1]{fontenc}
\usepackage{lmodern}
\providecommand{\notelanguage}{english}
\usepackage[finnish,english]{babel}
\AtBeginDocument{\expandafter\selectlanguage\expandafter{\notelanguage}}

\newcommand{\notesfinnishlanguage}{finnish}
\ifx\notelanguage\notesfinnishlanguage
\newcommand{\notestheoremname}{Lause}
\newcommand{\noteslemmaname}{Lemma}
\newcommand{\notespropositionname}{Propositio}
\newcommand{\notescorollaryname}{Korollaari}
\newcommand{\notesdefinitionname}{Määritelmä}
\newcommand{\notesexamplename}{Esimerkki}
\newcommand{\notesproblemname}{Tehtävä}
\newcommand{\notesexercisename}{Tehtävä}
\newcommand{\notesremarkname}{Huomautus}
\newcommand{\notessolutionname}{Ratkaisu}
\else
\newcommand{\notestheoremname}{Theorem}
\newcommand{\noteslemmaname}{Lemma}
\newcommand{\notespropositionname}{Proposition}
\newcommand{\notescorollaryname}{Corollary}
\newcommand{\notesdefinitionname}{Definition}
\newcommand{\notesexamplename}{Example}
\newcommand{\notesproblemname}{Problem}
\newcommand{\notesexercisename}{Exercise}
\newcommand{\notesremarkname}{Remark}
\newcommand{\notessolutionname}{Solution}
\fi

\usepackage[a4paper,margin=30mm]{geometry}
\usepackage{microtype}
\usepackage{graphicx}
\usepackage{enumitem}

\newcommand{\exerciseheading}[2]{%
  \par\noindent
  {\large\bfseries #1\par}
  \vspace{0.1em}
  \noindent{\large #2\par}
  \vspace{1.5em}
}

\usepackage{amsmath}
\usepackage{amssymb}
\usepackage{amsthm}
\usepackage{mathtools}
\usepackage{aliascnt}

% Row operations beside the affected row, e.g. \rowop{+2R_1}.
\newcommand{\rowop}[1]{\mathrel{\big|}\,#1}

% Linear systems: the argument gives the number of alignment pairs.
\newenvironment{linearsystem}[1]
{\left\{
  \begin{alignedat}{#1}}
    {
  \end{alignedat}\right.}

  \usepackage[hidelinks,pdfusetitle]{hyperref}

  \numberwithin{equation}{section}
  \setlist{itemsep=0.25em,topsep=0.5em}

  \theoremstyle{plain}
  \newtheorem{theorem}{\notestheoremname}[section]
  \newaliascnt{lemma}{theorem}
  \newtheorem{lemma}[lemma]{\noteslemmaname}
  \aliascntresetthe{lemma}
  \newaliascnt{proposition}{theorem}
  \newtheorem{proposition}[proposition]{\notespropositionname}
  \aliascntresetthe{proposition}
  \newaliascnt{corollary}{theorem}
  \newtheorem{corollary}[corollary]{\notescorollaryname}
  \aliascntresetthe{corollary}

  \theoremstyle{definition}
  \newaliascnt{definition}{theorem}
  \newtheorem{definition}[definition]{\notesdefinitionname}
  \aliascntresetthe{definition}
  \newaliascnt{example}{theorem}
  \newtheorem{example}[example]{\notesexamplename}
  \aliascntresetthe{example}
  \newaliascnt{problem}{theorem}
  \newtheorem{problem}[problem]{\notesproblemname}
  \aliascntresetthe{problem}
  \newtheorem{exercise}{\notesexercisename}

  \theoremstyle{remark}
  \newaliascnt{remark}{theorem}
  \newtheorem{remark}[remark]{\notesremarkname}
  \aliascntresetthe{remark}

  \newenvironment{solution}
  {
  \begin{proof}[\notessolutionname]}
    {
  \end{proof}}

  \usepackage[nameinlink,noabbrev]{cleveref}

  \crefname{theorem}{theorem}{theorems}
  \crefname{lemma}{lemma}{lemmas}
  \crefname{proposition}{proposition}{propositions}
  \crefname{corollary}{corollary}{corollaries}
  \crefname{definition}{definition}{definitions}
  \crefname{example}{example}{examples}
  \crefname{problem}{problem}{problems}
  \crefname{exercise}{exercise}{exercises}
  \crefname{remark}{remark}{remarks}

  \ifx\notelanguage\notesfinnishlanguage
  \crefname{theorem}{lause}{lauseet}
  \crefname{lemma}{lemma}{lemmat}
  \crefname{proposition}{propositio}{propositiot}
  \crefname{corollary}{korollaari}{korollaarit}
  \crefname{definition}{määritelmä}{määritelmät}
  \crefname{example}{esimerkki}{esimerkit}
  \crefname{problem}{tehtävä}{tehtävät}
  \crefname{exercise}{tehtävä}{tehtävät}
  \crefname{remark}{huomautus}{huomautukset}
  \fi

  \newcommand{\NN}{\mathbb{N}}
  \newcommand{\NNone}{\mathbb{N}_1}
  \newcommand{\ZZ}{\mathbb{Z}}
  \newcommand{\QQ}{\mathbb{Q}}
  \newcommand{\RR}{\mathbb{R}}
  \newcommand{\CC}{\mathbb{C}}
  \newcommand{\PP}{\mathbb{P}}

  \DeclareMathOperator{\im}{im}
  \DeclareMathOperator{\rank}{rank}
  \DeclareMathOperator{\syt}{syt}
  \DeclareMathOperator{\pyj}{pyj}
  \DeclareMathOperator{\dom}{dom}
  \DeclarePairedDelimiter{\abs}{\lvert}{\rvert}
  \DeclarePairedDelimiter{\norm}{\lVert}{\rVert}
  \newcommand{\set}[1]{\left\lbrace#1\right\rbrace}
  \newcommand{\maxset}[1]{\max\set{#1}}
  \newcommand{\minset}[1]{\min\set{#1}}
  \newcommand{\powerset}[1]{\mathcal{P}\!\left(#1\right)}


\usepackage{tikz}

\title{Analyysi II\\Ryhmätyö 3 ratkaisut}
\author{}
\date{}

\begin{document}

\exerciseheading{Analyysi II}{Ryhmätyö 3 ratkaisut}

\begin{exercise}[Protter--Morrey, lause 13.2 reaaliluvuilla]
  Käy huolellisesti läpi kiintopistelauseen 13.2 todistus ja kirjoita
  siitä mahdollisimman selkeä esitys avaruudessa \(\RR\), jossa etäisyys
  on \(d(x,y)=\lvert x-y\rvert\). Lisää tarpeelliset selitykset;
  voit myös muokata kirjan todistusta. Todistettava väite on seuraava:

  Olkoon \(f\colon\RR\to\RR\) kuvaus ja olkoon \(k\in(0,1)\) siten, että
  \[
    \lvert f(x)-f(y)\rvert\leq k\lvert x-y\rvert
    \qquad\text{kaikilla }x,y\in\RR.
  \]
  Tällöin on olemassa täsmälleen yksi \(x_*\in\RR\), jolle
  \(f(x_*)=x_*\). Lisäksi jokaisella alkuarvolla \(x_1\in\RR\) jono,
  joka määritellään asettamalla
  \[
    x_{n+1}=f(x_n),\qquad n=1,2,\ldots,
  \]
  suppenee pisteeseen \(x_*\).
\end{exercise}

\begin{exercise}[Protter--Morrey, luku 13.1, tehtävä 8]
  Olkoon \(I=[0,\infty)\) ja \(f\colon I\to I\),
  \[
    f(x)=\sqrt{1+x^2}.
  \]
  Osoita, että
  \[
    \lvert f(x)-f(y)\rvert<\lvert x-y\rvert
    \qquad\text{kaikilla }x,y\in I,\ x\neq y,
  \]
  mutta funktiolla \(f\) ei ole kiintopistettä välillä \(I\).
  Päättele, ettei lause 13.2 päde, jos sen oletus \(0<k<1\)
  korvataan ehdolla \(k=1\).
\end{exercise}

\begin{exercise}[Protter--Morrey, luku 13.1, tehtävä 9]
  Todista lauseen 13.2 seuraus seuraavassa täsmällisessä muodossa:

  Olkoot \(a,b\in\RR\), \(a<b\), ja \(I=[a,b]\).
  Olkoon \(f\colon I\to I\) jatkuva välillä \(I\) ja derivoituva
  välillä \((a,b)\). Oletetaan, että on olemassa vakio \(k\in(0,1)\),
  jolle
  \[
    \lvert f'(x)\rvert\leq k
    \qquad\text{kaikilla }x\in(a,b).
  \]
  Tällöin on olemassa täsmälleen yksi \(x_*\in I\), jolle
  \(f(x_*)=x_*\).
\end{exercise}

\begin{exercise}[Protter--Morrey, luku 13.1, tehtävä 10]
  Olkoon \(f\colon\RR^2\to\RR^2\) kuvaus, joka määritellään asettamalla
  \[
    f(x,y)=\left(\frac14x+\frac13y-2,\;
    \frac15x-\frac13y+3\right).
  \]
  Osoita lauseen 13.2 avulla, että kuvauksella \(f\) on kiintopiste,
  eli on olemassa \((x_*,y_*)\in\RR^2\), jolle
  \(f(x_*,y_*)=(x_*,y_*)\).
\end{exercise}

\end{document}
